\scanpage{97}

On conclut de (7) que $\ell^N$ est \uline{complet}. Il suffit de prouver que \deleted{toute suite} $(x^k)$ \deleted{dans $\ell^N$} d'éléments $x^k\in\ell^N$, telle que $\sum\lVert x^{k}\rVert_N<+\infty$, est sommable dans $\ell^N$. Or \deleted{elle} on aura à fortiori $\sum\lVert x^k\rVert_\infty<+\infty$, donc la \deleted{série} suite sera sommable dans $\ell^\infty$, \inserted{vers un $x\in\ell^\infty$,} d'autre part l'on \deleted{aura} voit sur (7) que $x\in\ell^N$. D'autre part,

% slanted
\margin{Inutile, il suffit de procéder par dualité\,: $\ell^N$ est le dual de $\ell_N$.}

on a

\[
N\left(x-\sum_{k=1}^n x^k\right)=\deleted{\unsure{\leq}}N\left(\deleted{\unsure{x-|x}}\sum_{n+1}^\infty x^k\right)\leq\sum_{n+1}^{\infty}N(x^k),
\]

qui tend vers \deleted{\unsure{l'$\infty$} \unsure{puisque la|pour}}\inserted{$0$ pour $n\to+\infty$}, puisque $\sum N(x^k)$ est une série convergente. Il s'ensuit que l'on a aussi $\sum x^k=x$ \uline{dans $\ell^N$}.

Soit $\ell_N$ l'adhérence de $C^{(\underline N)}$ dans $\ell^N$, comme $\ell^N$ est complet, $\ell_N$ s'identifie au complété de $C^{(\underline N)}$ pour la norme $N$. Il est encore trivial que si $N'\leq N''$, alors $\ell_{N'}\subset\ell_{N''}$, et l'application identique $\ell_{N'}\to\ell_{N''}$ est de norme $\leq 1$. En particulier, les espaces $\ell_N$ associés à $N_1$ et $N_\infty$ étant resp. $\ell^1$ et $c_0$, on obtient

\[
\ell^1\subset\ell_N\subset c_0\tag*{(8)}
\]

\uline{Proposition} On a $\ell^N\subset c_0$ sauf si $N$ est équivalente à $N_\infty$:

\deleted{Notons la formule, qui précise la \unsure{condition|propriétés} \unsure{ci-dessus|suivantes}}

\[
\deleted{\underline{\ell}_N=\underline{\ell}^N\cap\underline{c}_0}\tag*{\deleted{(9)}}
\]

\deleted{Il suffit de prouver que $x\in\ell^N\cap c_0$ implique $x\in\ell_N$. \unsure{Et} En effet, on aura}

Notons que $\deleted{\unsure{C}}R^{(N)}$ est en dualité avec lui-$=$ par l'accouplement $\langle x,y\rangle=\sum x_i y_i$. Cela permet de faire correspondre à \deleted{toute} une norme $N$ sur $R^{(N)}$, la \uline{norme polaire}, notée $N'$, définie par

\[
N^\circ(y')=\underset{N(x)\leq 1}{\operatorname{Sup}}\,|\langle x,y'\rangle|\tag*{(9)}
\]

% margin note written sideways
\margin{De plus, \deleted{\unsure{comme|clairement}} comme $(N^\circ)^\circ=N$, on peut le vérifier .... \deleted{\unsure{on est ramené} \illegible{2}}}

Si $N$ est une norme de Schatten, on vérifie trivialement que $N^\circ$ est une norme de Schatten, \deleted{et} \unsure{On dira que} \deleted{le dual de $\ell_N$ s'identifie à $\ell^{N'}$, l'accouplement étant $\sum x_i x_i'$ est \unsure{abst} sommable, et que} \inserted{si $u\in\ell_N$, $u'\in\ell^{N^\circ}$, alors \deleted{\unsure{on a le}}}
% archive sheet number: 98
